\section{Discussion}
\label{sec:discussion}

\Cref{tab:headline} collects the measurements that carry our reading.
The same residual algebra $x+f(x)$ supports two different transports, and refinement
at fixed horizon makes the difference visible without retraining either model.

\begin{table}[t]
    \centering
    \caption{Headline measurements on ImageNet-1k.
    Top-1 and ODE limit from \cref{tab:r-ablation}; straightness (Str.) and consecutive
    block cosine at $R{=}1$ (\cref{tab:coherence,app:arch-transport}); parameters from
    \cref{sec:weight_tied}.
    Swin-T and the ResNets are torchvision checkpoints that we did not retrain.}
    \label{tab:headline}
    \small
    \setlength{\tabcolsep}{2.6pt}
    \begin{tabular}{lrrrrrr}
        \toprule
        Model & Params & $R{=}1$ & ODE & Gap & Str. & $\cos_{k,k+1}$ \\
        \midrule
        ConvNeXt-T         & 28.59M & 80.96 & 80.23 & $-0.73$ & 0.50 & 0.16 \\
        Weight-tied S3     & 18.97M & 80.34 & 80.12 & $-0.22$ & 0.72 & 0.77 \\
        All-stages tied    & 8.68M  & 78.99 & --    & --      & --   & --   \\
        Swin-T             & --     & 81.08 & 80.32 & $-0.76$ & 0.80 & 0.43 \\
        ResNet-50          & --     & 80.34 & 79.43 & $-0.91$ & 0.37 & $-0.11$ \\
        ResNet-101         & --     & 81.67 & 81.05 & $-0.62$ & 0.16 & 0.01 \\
        \bottomrule
    \end{tabular}
\end{table}

\paragraph{A staircase path.}
The untied ConvNeXt-T stage does not move back and forth: every block's increment survives to the end of the stage ($s_k>0$ for every image and
block), and consecutive increments are positively correlated (cosine $0.17$) where
random blocks of matched strength are uncorrelated.
Each block writes mostly in a new direction, so the net displacement grows more like
the square root of the number of steps than like their sum.
Straightness places it halfway between random blocks and the shared field
($0.33$, $0.50$, $0.78$), and the detour is made at the block switches, where the
heading turns by $80^\circ$ and $45\%$ of the coordinates of the velocity change sign;
inside a block the path is a shallow arc whose length exceeds its chord by $1.5\%$.
These statistics are already present at native resolution (for the untied stage
$R{=}1$ and $R{=}100$ agree to within $0.01$ in straightness and cosine), so they describe the trained
network and not the refinement.

\paragraph{Why does one field nearly match nine?}
The shared stage reaches $80.34\%$ against $80.96\%$ with $33\%$ fewer parameters.
Our measurements are consistent with a simple reading: the shared field, which must act
sensibly on every state it visits, moves the state along a coherent direction and
reaches the decoder's operating region earlier (prefix top-1 $74.7\%$ at $t{=}3$
against $30.6\%$), while the nine untied blocks spread the same transport over nine
nearly orthogonal steps and make most of the final displacement in the last blocks.
They do not show that the untied directions are unnecessary, nor what the $0.62$ points
buy.
Nearly orthogonal increments could carry information that a single field cannot
write, and the accuracy gap itself comes from runs with different learning rates and
warmups (\cref{app:training}).
Why the coherent field nearly matches the staircase remains an open question.

\paragraph{Refinement as a probe.}
Every architecture converges to a continuous limit at fixed $T$ that keeps the function
within $0.22$ to $0.91$ points, although it sits $11^\circ$ to $21^\circ$ from the native
state (\cref{app:solver-endpoints}); breaking the horizon instead collapses every
network.
A trained stage is thus a sequence of one-step solves, each replaceable by a finer
integral of the same block, but not a Neural ODE with weights smooth in depth
(\cref{app:bilinear}).
The probe needs no retraining, the Gram statistics only need block increments and prefix
decoding only needs the trained downstream network, so all three apply to any residual
stream, including looped transformers~\citep{dehghani2019universal,giannou2023looped},
the weight-tied counterpart of our shared stage.

\paragraph{Limitations.}
All results are on ImageNet-1k classification, and the trained tied and untied stages
are one pair of ConvNeXt-T runs, without seeds and with different learning-rate
schedules.
The open question above is the main one: we describe how the two transports differ,
not why the difference costs only $0.62$ points.
The channel basis is not privileged, so per-coordinate readings of the trajectories
(sign changes between blocks) are partly determined by the Gram matrix and partly by the
basis (\cref{app:per-coordinate}); we base our claims on Gram statistics only.
The metrics exclude the massive ConvNeXt channel and the Swin corner tokens, and we
have not tested whether the ReLU after the ResNet sum explains its negatively correlated
increments (\cref{app:arch-transport}).

\paragraph{Toward more expressivity per parameter.}
The two regimes suggest a design space between them: a shared field with small
per-block corrections, or cycles of a few shared blocks, would interpolate between an
autonomous and a switched flow.
Related schemes exist for convolutional networks~\citep{savarese2019learning} and
transformers, through cyclic layer sharing~\citep{takase2021lessons} and shared layers
with low-rank updates~\citep{bae2024relaxed}.
The coherence statistics place such models on the ladder from random to autonomous
and can test whether coherence and accuracy move together.
